\begin{lema}[Regularidad del perfil en el extremo]\label{lem:regularidad-extremo}
Para algún \(x_0>0\), \(G\) es analítica en un entorno de \([0,x_0]\).
En ese entorno \(\numerr\) es meromorfa con único polo posible en cero, de orden a lo sumo cuatro.
\end{lema}
\begin{proof}
El perfil continuo de masa uno tiene \(\nu_1\ge c>0\), marginales axiales uno y
\(u_f\ge c>0\) en ambas familias oblicuas~\cite{pv}. Reducimos \(c\) para que \(c<1\).
En un disco complejo \(|x|<r\) se tiene \(|U_f|\ge c/2\), luego \(W\) y \(1/U_f\)
son holomorfas como funciones acotadas de la posición. Todos los dominios de cruces, fibras,
paridades e incidencias de la sección~\ref{sec:contracciones} son independientes de \(x\).

\paragraph{La inversa y sus entradas.}
Usemos una medida fija \(d\sigma=dy\) en cada eje y \(dy/2\) en cada hoja oblicua. Sean
\(D_x(a)=1\) en los ejes y \(D_x(a)=x\) en las oblicuas, y \(u_x(a)=1\) o \(U_a(x)\),
respectivamente. Así \(d\rho=(u_x/D_x)d\sigma\). Para dos líneas distintas que se cortan,
\(K(a,b)=\nu_1(a\cap b)\eta_a\eta_b\), con \(\eta_a=D_x(a)/u_x(a)\).
Por tanto \(K=D_x K_0 D_x\), donde el núcleo \(K_0\) es holomorfo y uniformemente
acotado, incluso en los bordes al usar la convención de cruce fijada. Las entradas restantes
valen cero. En \(L^2(d\sigma)\), los productos de núcleos usan \(d\sigma\) y los factores
\(D_x,u_x\) actúan por multiplicación. El operador que se invierte es
\[
 \mathcal T_x=I+u_x K_0D_x+\Pi .
\]
En \(x=0\) es triangular por bloques (ejes primero), con identidad en el bloque oblicuo.
El bloque axial es la identidad más el operador de bloques cruzados \(\nu_1,\nu_1^{\mathsf T}\)
más \(\Pi\). En las constantes axiales común y alternada sus autovalores son \(2,1\).
En las funciones de media cero por eje la norma del operador cruzado es \(\le1-c\):
el núcleo \((\nu_1-c)/(1-c)\) es bistocástico, luego es una contracción en \(L^2\)
por Jensen e integración. Esto da la inversa acotada del bloque axial por serie geométrica.
El otro bloque fuera de la diagonal es acotado por tener núcleo acotado y medida finita.
La inversa triangular existe. Reducimos \(r\) para que
\(\|\mathcal T_0^{-1}(\mathcal T_x-\mathcal T_0)\|<1/2\). La serie de Neumann da
una inversa holomorfa uniformemente acotada en ese disco.

La parte de \(L_0+3\Pi\) con extremos \(a,b\) vale
\[
 D_x(a)D_x(b)\int K_0(a,t)D_x(t)
     [\mathcal T_x^{-1}(u_x K_0(\cdot,b))](t)\,d\sigma(t).
\]
Las dos secciones exteriores tienen normas \(L^2\) uniformes para \emph{cada} extremo,
por la fórmula de cruce. La integral bilineal es holomorfa y acotada también cuando \(a=b\).
Así se obtiene la versión del núcleo que puede evaluarse en una casilla marcada o sobre la
diagonal. En \(\theta_a\), el factor \(K(a,t)^2\) deja \(D_x(t)^2\), que absorbe el único
\(1/D_x(t)\) de la medida. En \(s_a\) y en la última suma de \(\mathcal V\) las fibras y la
medida de conteo normalizada de las líneas son fijas. Las fórmulas integrales del
apéndice~\ref{app:predictor} dan por tanto \(\mathcal V\) holomorfa.

\paragraph{El funcional de torres.}
En \(x=0\), el escalado de \(v_0=u_Du_A\) es \(\nu_1\) por su factorización.
En \(C[0,1]\) la ecuación para el logaritmo \(z\) del factor simétrico es
\(e^{z(r)}\int v_x(r,q)e^{z(q)}\,dq=1\). Su derivada en la solución inicial es
\(I+\mathcal P\), donde \(\mathcal P\) tiene núcleo \(\nu_1\). Vale dos sobre las constantes.
En media cero, \(\|\mathcal P\|_\infty\le1-c\) por la misma descomposición bistocástica.
El teorema de la función implícita en este espacio complejo da un escalado holomorfo
en un disco menor. Su rama real coincide con el escalado positivo por unicidad.
Para el núcleo centrado \(\centrado=p_x-1\) del apéndice~\ref{app:predictor},
\(\|\centrado(0)\|_{2\to2}\le1-c\). Reduciendo el disco, \((I-\centrado^2)^{-1}\) es holomorfa acotada.
El núcleo de \(T=\centrado(I-\centrado^2)^{-1}\centrado\) se evalúa entre dos secciones acotadas de
 \(\centrado\),
incluso en la diagonal, como en el párrafo anterior. Los productos de núcleos y
\(\operatorname{tr}(T^2)=\iint T(r,q)T(q,r)\,dr\,dq\) preservan la holomorfía.
Así \(F\), y con ella \(G\), son holomorfas cerca de cero.

\paragraph{Orden máximo en las contracciones.}
En un producto directo, la esperanza de los indicadores es el producto de los \(U_\ell\)
de las líneas distintas sin marcar. Es holomorfo incluso para dos cruces que comparten una
transversal. Se usa este producto antes de cualquier división por la disponibilidad común.
Las medias de área usan \(\nu_1\,di\), y la fibra de una casilla marcada usa
\(\eta_\lambda\nu_1\,dy\), ambas acotadas. Cada integral libre \(d\rho\) cuesta a lo sumo
un factor \(x^{-1}\). La tabla cuenta esos factores en las expresiones físicas del
apéndice~\ref{app:interior}, con los cuadrados y sus dos correcciones donde corresponden.
\begin{center}\small
\begin{tabular}{@{}lp{.54\linewidth}@{}}\toprule
Contracción & Fuentes posibles de polos\\ \midrule
\(G_{\rm cof}, f2E, f2F,f3F,f3c,f4,f5c,f5F,f7\) & Sólo áreas, centros y fibras marcadas.\\
\(f2z\) & Una línea libre y el multiplicador acotado \(u_f/U_f\).\\
\(f6\) & \(\Lambda,\omega\) usan cruces y fibras. La respuesta \(q\) integra una línea libre.\\
\(C_{3,\rm cuad}\) & La potencia explícita \(x^{n_{RC}-3}\) cuesta a lo sumo tres factores.\\
\(Q_1+\Delta_1\) & Una copia de \(\mathcal E_2\) integra a lo sumo una línea libre.\\
\(Q_2+\Delta_2\) & Dos líneas exteriores y a lo sumo una interior en cada copia de
 \(\mathcal E_2\).\\ \bottomrule
\end{tabular}
\end{center}
En las dos últimas filas se usan \(L_1,L_2^{\rm int}\), antes de descomponer los
productos en covarianzas. La suma discreta con \(\mathcal S(w,u)\) recorre las
transversales de los cruces \(u\cap j\), un número acotado. Cada producto en la tabla,
multiplicado por \(x^4\), queda así holomorfo bajo medidas fijas, con mayorante constante
integrable en un disco menor. Integrar la serie por Cauchy en un círculo interior
da la holomorfía de \(x^4d(x)\). Puesto que \(\numerr=d+G-xG'\), resulta el enunciado.
\end{proof}
